% GENERATED FILE. Edit canonical lessons, then run npm run generate:books.
% canonical-source-hash: fnv1a64:9bf6929d1374bffa
% canonical-lessons: AR-C204-nahwa, AR-C204-mintaqa, AR-C204-markaz, AR-C204-asima, AR-C204-ayna

\chapter{Towards, Area, Centre, Capital city, Where}
\label{ch:ar-towards-area-centre-capital-city-where}
\hlchaptermodality{204}

\begin{chapteropening}
\textbf{By the end of this chapter:} \emph{I can say towards, an area, a centre, a capital city and where.}

Complete the last lesson of chapter 204: I can say towards, an area, a centre, a capital city and where.
\end{chapteropening}

\section[\texorpdfstring{\ar{نحو} (naḥwa)}{نحو (naḥwa)}]{\ar{نحو} (naḥwa) — towards}

\label{lesson:AR-C204-nahwa}



\begin{quote}
Before the new one: say the Arabic for a century, then the Arabic for weekly.
\end{quote}

\subsection*{You'll want to know: \ar{نحو}}
\textbf{\ar{نحو}} — \emph{naḥwa} — \textquotedblleft{}towards\textquotedblright{}.

\begin{cousinweb}
\textbf{\ar{نحو}} (\emph{naḥwa}) is built on the native Arabic root \emph{n-ḥ-w}, \textquotedblleft{}to go towards\textquotedblright{} (Lane). \emph{Naḥwa} is \textquotedblleft{}towards\textquotedblright{}; the same word names Arabic grammar, \emph{naḥw}.
\end{cousinweb}

\begin{grammarlens}[title={What you've built}]
One more word for this chapter.
\end{grammarlens}

\subsection*{Your turn}
\begin{itemize}
\raggedright
  \item \emph{Say it:} \emph{naḥwa}
  \item \emph{Say it:} \emph{naḥwa}, once more
  \item \emph{Say it:} say \emph{naḥwa}
  \item \emph{Recall:} say the Arabic for a century, then the Arabic for weekly, then say \emph{naḥwa} again
\end{itemize}

\begin{tcolorbox}[breakable,colback=teal!4,colframe=teal!35!black,title={Before you move on}]
What does \ar{نحو} mean? (\textquotedblleft{}Towards\textquotedblright{}.) Say it once more.
\end{tcolorbox}
\section[\texorpdfstring{\ar{منطقة} (minṭaqa)}{منطقة (minṭaqa)}]{\ar{منطقة} (minṭaqa) — an area, a district}

\label{lesson:AR-C204-mintaqa}



\begin{quote}
Before the new one: say the Arabic for weekly, then the Arabic for towards.
\end{quote}

\subsection*{You'll want to know: \ar{منطقة}}
\textbf{\ar{منطقة}} — \emph{minṭaqa} — \textquotedblleft{}an area, a district\textquotedblright{}.

\begin{cousinweb}
\textbf{\ar{منطقة}} (\emph{minṭaqa}) is built on the native Arabic root \emph{n-ṭ-q}, \textquotedblleft{}to speak; a girdle\textquotedblright{} (Lane). Lane's \emph{minṭaqa} is a girdle; the zone, like Greek \emph{zōnē}, \textquotedblleft{}belt\textquotedblright{}, grew from it.
\end{cousinweb}

\begin{grammarlens}[title={What you've built}]
One more word for this chapter.
\end{grammarlens}

\subsection*{Your turn}
\begin{itemize}
\raggedright
  \item \emph{Say it:} \emph{minṭaqa}
  \item \emph{Say it:} \emph{minṭaqa}, once more
  \item \emph{Say it:} say \emph{minṭaqa}
  \item \emph{Recall:} say the Arabic for weekly, then the Arabic for towards, then say \emph{minṭaqa} again
\end{itemize}

\begin{tcolorbox}[breakable,colback=teal!4,colframe=teal!35!black,title={Before you move on}]
What does \ar{منطقة} mean? (\textquotedblleft{}An area, a district\textquotedblright{}.) Say it once more.
\end{tcolorbox}
\section[\texorpdfstring{\ar{مركز} (markaz)}{مركز (markaz)}]{\ar{مركز} (markaz) — a centre}

\label{lesson:AR-C204-markaz}



\begin{quote}
Before the new one: say the Arabic for towards, then the Arabic for an area.
\end{quote}

\subsection*{You'll want to know: \ar{مركز}}
\textbf{\ar{مركز}} — \emph{markaz} — \textquotedblleft{}a centre\textquotedblright{}.

\begin{cousinweb}
\textbf{\ar{مركز}} (\emph{markaz}) is built on the native Arabic root \emph{r-k-z}, \textquotedblleft{}to plant (a spear) upright\textquotedblright{} (Lane). Lane's \emph{markaz} is the place where a spear is planted upright; the centre grew from it.
\end{cousinweb}

\begin{grammarlens}[title={What you've built}]
One more word for this chapter.
\end{grammarlens}

\subsection*{Your turn}
\begin{itemize}
\raggedright
  \item \emph{Say it:} \emph{markaz}
  \item \emph{Say it:} \emph{markaz}, once more
  \item \emph{Say it:} say \emph{markaz}
  \item \emph{Recall:} say the Arabic for towards, then the Arabic for an area, then say \emph{markaz} again
\end{itemize}

\begin{tcolorbox}[breakable,colback=teal!4,colframe=teal!35!black,title={Before you move on}]
What does \ar{مركز} mean? (\textquotedblleft{}A centre\textquotedblright{}.) Say it once more.
\end{tcolorbox}
\section[\texorpdfstring{\ar{عاصمة} (ʿāṣima)}{عاصمة (ʿāṣima)}]{\ar{عاصمة} (ʿāṣima) — a capital city}

\label{lesson:AR-C204-asima}



\begin{quote}
Before the new one: say the Arabic for an area, then the Arabic for a centre.
\end{quote}

\subsection*{You'll want to know: \ar{عاصمة}}
\textbf{\ar{عاصمة}} — \emph{ʿāṣima} — \textquotedblleft{}a capital city\textquotedblright{}.

\begin{cousinweb}
\textbf{\ar{عاصمة}} (\emph{ʿāṣima}) is built on the native Arabic root \emph{ʿ-ṣ-m}, \textquotedblleft{}to protect, to hold back\textquotedblright{} (Lane). \emph{ʿĀṣima} is the feminine active participle, \textquotedblleft{}one that protects\textquotedblright{}; the capital city is a modern sense.
\end{cousinweb}

\begin{grammarlens}[title={What you've built}]
One more word for this chapter.
\end{grammarlens}

\subsection*{Your turn}
\begin{itemize}
\raggedright
  \item \emph{Say it:} \emph{ʿāṣima}
  \item \emph{Say it:} \emph{ʿāṣima}, once more
  \item \emph{Say it:} say \emph{ʿāṣima}
  \item \emph{Recall:} say the Arabic for an area, then the Arabic for a centre, then say \emph{ʿāṣima} again
\end{itemize}

\begin{tcolorbox}[breakable,colback=teal!4,colframe=teal!35!black,title={Before you move on}]
What does \ar{عاصمة} mean? (\textquotedblleft{}A capital city\textquotedblright{}.) Say it once more.
\end{tcolorbox}
\section[\texorpdfstring{\ar{أين} (ayna)}{أين (ayna)}]{\ar{أين} (ayna) — where}

\label{lesson:AR-C204-ayna}



\begin{quote}
Before the new one: say the Arabic for a centre, then the Arabic for a capital city.
\end{quote}

\subsection*{You'll want to know: \ar{أين}}
\textbf{\ar{أين}} — \emph{ayna} — \textquotedblleft{}where\textquotedblright{}.

\begin{cousinweb}
\textbf{\ar{أين}} (\emph{ayna}) is listed by Lane under \emph{ʾ-y-n} as an adverbial noun that asks about place. It is not built on a verbal root; Hebrew \emph{ʾayyē}, \textquotedblleft{}where\textquotedblright{}, is usually compared.
\end{cousinweb}

\begin{grammarlens}[title={What you've built}]
One more word for this chapter.
\end{grammarlens}

\subsection*{Your turn}
\begin{itemize}
\raggedright
  \item \emph{Say it:} \emph{ayna}
  \item \emph{Say it:} \emph{ayna}, once more
  \item \emph{Say it:} say \emph{ayna}
  \item \emph{Recall:} say the Arabic for a centre, then the Arabic for a capital city, then say \emph{ayna} again
\end{itemize}

\begin{tcolorbox}[breakable,colback=teal!4,colframe=teal!35!black,title={Before you move on}]
What does \ar{أين} mean? (\textquotedblleft{}Where\textquotedblright{}.) Say it once more.
\end{tcolorbox}
